The within-group discrete ordinates equations within a uniform material can be written as:
\begin{equation}
  \label{eq:sn-eqs}
  \underbrace{\ohat_{x,n}\pdv{\psing}{x}+\ohat_{y,n}\pdv{\psing}{y}+\ohat_{z,n}\pdv{\psing}{z}}_{\ohat_n\cdot\nabla\psing \rarg}+\Sigma_t^g \psing\rarg=q_n^g\rarg,\quad n\in \{1,2,\ldots,N\}.
\end{equation}
These equations are linked together via a quadrature set for the angle-integrated flux:
\begin{equation}
  \phi^g\rarg := \int_{4\pi}\psi^g\ro\odif{\ohat}\approx \sum_{n=1}^{N}w_n\psing \rarg.
\end{equation}
The source term in \cref{eq:sn-eqs} will be assumed to be comprised of scattering, fission, and an independent source:
\begin{equation}
  q_n^g\rarg:= \sum_{g'=1}^G \int_{4\pi}\Sigma_s^{g'\rightarrow g}\left(\rvec,\ohat'\cdot\ohat\right)\psi^{g'}\left(\rvec,\ohat'\right)\odif{\ohat'}+\frac{\chi^g}{4\pi \keff } \sum_{g'=1}^G  \nu^{g'}\Sigma_f^{g'}\phi^{g'}\rarg+q^g_{\text{ind}}\left(\rvec,\ohat_n\right).
\end{equation}
Some important notes:
\begin{itemize}
  \item The double-difference scattering cross section within the scattering source is assumed to be only a function of scattering angle, not the incoming angle.
  \item Delayed neutrons are not considered.
\end{itemize}
Going forward, we drop the $n$ and $g$ indices where it is obvious they are implied.

First, we partition the spatial domain into $E$ non-overlapping subdomains of unstructured hexahedral elements.
We then require \cref{eq:sn-eqs} to hold over each element.
Multiplying \cref{eq:sn-eqs} by a test function, $v(x,y,z)$, integrating over all space, and applying vector identities we arrive at the weak form of the \gls{fobte}:
\begin{equation}
  \label{eq:weak-form}
  \oint v\psi\ohat\cdot\hat{n}\odif{S}-\int \psi\ohat\cdot\nabla v\odif{V}+\int\Sigma_t\psi v\odif{V}=\int qv\odif{V}
\end{equation}

Next, we approximate the discrete-ordinate, within-group angular flux and source with a spectral expansion using Lagrange polynomials:
\begin{subequations}
  \label{eq:spectral-expansion}
  \begin{equation}
    \psing\rarg\approx \sum_{i=0}^{P}\sum_{j=0}^{P}\sum_{k=0}^{P}\psi_{ijk}\ell_i\left(x\right)\ell_j\left(y\right)\ell_k\left(z\right)
  \end{equation}
  \begin{equation}
    q_n^g\rarg \approx \sum_{i=0}^{P}\sum_{j=0}^{P}\sum_{k=0}^{P}q_{ijk}\ell_i\left(x\right)\ell_j\left(y\right)\ell_k\left(z\right)
  \end{equation}
\end{subequations}
where the nodes of the Lagrange interpolants are selected to be the mapped nodes belonging to the $P$\textsuperscript{th}-order \gls{gll} quadrature set.
From \cite{devilleHighorderMethodsIncompressible2002}:
\begin{equation}
  \label{eq:gll}
  \begin{aligned}
     & \int_{-1}^1 u(\xi)\odif{\xi} & \approx & \sum_{k=0}^{P}\rho_k u\left(\xi_k\right)                                  \\
     & \xi_k                        & =       & \begin{cases}
                                                  -1,                        & k=0              \\
                                                  \text{zeros of }L'_P(\xi), & 1\leq k \leq P-1 \\
                                                  1,                         & k=P              \\
                                                \end{cases}                           \\
     & \rho_k                       & =       & \frac{2}{P\left(P+1\right)}\frac{1}{\left[L_P\left(\xi_k\right)\right]^2}
  \end{aligned}
\end{equation}
The above rule is exact for polynomials of degree up to $2P-1$.
Note that for $\rvec \in \mathbb{R}^n$, for $n>1$, a single quadrature set is created based on \cref{eq:gll}, then points are mapped to the higher space via a tensor product.

To form the spectral element model, we insert \cref{eq:spectral-expansion} into \cref{eq:weak-form}.
Additionally, we choose $v(x,y,z)$ to belong to the same space as $\psi(x,y,z)$, forming a Galerkin projection:
\begin{equation}
  v(x,y,z)=\sum_{\hat{\imath}=0}^{P}\sum_{\hat{\jmath}=0}^{P}\sum_{\hat{k}=0}^{P}v_{\hat{\imath}\hat{\jmath}\hat{k}}\ell_{\hat{\imath}}(x)\ell_{\hat{\jmath}}(y)\ell_{\hat{k}}(z).
\end{equation}
\subsubsection{Curvilinear Elements}
To represent curved geometry, such as fuel pins and coolant channels, each hexahedral element is mapped from the reference cube $(\mu,\eta,\xi)\in[-1,1]^3$ \emph{isoparametrically}, i.e., using the same Lagrange basis as the solution \cite{devilleIterativeSolutionIsoparametric1987,devilleHighorderMethodsIncompressible2002}:
\begin{equation}
  \label{eq:isoparametric-map}
  \rvec(\mu,\eta,\xi)=\sum_{i=0}^{P}\sum_{j=0}^{P}\sum_{k=0}^{P}\rvec_{ijk}\,\ell_i(\mu)\ell_j(\eta)\ell_k(\xi),
\end{equation}
where $\rvec_{ijk}$ are the physical coordinates of the element's \gls{gll} nodes.
The element faces are therefore polynomial surfaces of degree $P$, and on each element the basis functions in \cref{eq:spectral-expansion} are understood as $\ell_i(\mu)\ell_j(\eta)\ell_k(\xi)$, with $(\mu,\eta,\xi)$ the reference coordinates of the point $\rvec$.
Writing $(r_1,r_2,r_3)=(\mu,\eta,\xi)$, the geometry enters through the Jacobian and the metric gradients,
\begin{equation}
  \label{eq:geometric-factors}
  J:=\det\left(\pdv{\rvec}{\mu,\eta,\xi}\right),\qquad
  \nabla\mu=\frac{1}{J}\pdv{\rvec}{\eta}\times\pdv{\rvec}{\xi},\quad
  \nabla\eta=\frac{1}{J}\pdv{\rvec}{\xi}\times\pdv{\rvec}{\mu},\quad
  \nabla\xi=\frac{1}{J}\pdv{\rvec}{\mu}\times\pdv{\rvec}{\eta},
\end{equation}
which vary within a curved element and are evaluated at each \gls{gll} node by differentiating \cref{eq:isoparametric-map}.
By the chain rule, the streaming derivative of the test function is
\begin{equation}
  \ohat\cdot\nabla v=\sum_{d=1}^{3}\left(\ohat\cdot\nabla r_d\right)\pdv{v}{r_d}.
\end{equation}

We evaluate every volume integral in \cref{eq:weak-form} with the tensor-product \gls{gll} rule, $\int_e f\odif{V}\approx\sum_m W_m J_m f(\rvec_m)$, where $m = i + j(P+1)+k(P+1)^2$ indexes the volume nodes (and $\hat{m} = \hat{\imath} + \hat{\jmath}(P+1)+\hat{k}(P+1)^2$ the test functions), $W_m:=\rho_i\rho_j\rho_k$, and a subscript $m$ denotes evaluation at node $m$.
Derivatives of the basis at the nodes require only the one-dimensional \gls{gll} differentiation matrix,
\begin{equation}
  \label{eq:gll-diff}
  \hat{D}_{p\hat{\imath}}:=\odv{\ell_{\hat{\imath}}}{\mu}\bigg|_{\mu_p},
\end{equation}
applied along one reference direction at a time:
\begin{equation}
  D_1:=I\otimes I\otimes\hat{D},\qquad D_2:=I\otimes\hat{D}\otimes I,\qquad D_3:=\hat{D}\otimes I\otimes I,
\end{equation}
where $I$ is the $(P+1)\times(P+1)$ identity.
Because the quadrature points coincide with the interpolation nodes, the mass matrix remains diagonal even on a curved element, and all of the geometry collects into diagonal matrices of nodal values:
\begin{subequations}
  \label{eq:sem-general-matrices}
  \begin{align}
    \label{eq:general-mass}
    M                     & \approx \operatorname{diag}\left(W_mJ_m\right),                                                                          \\
    \label{eq:general-streaming}
    S_{\ohat}             & \approx \sum_{d=1}^{3}D_d^{T}\Lambda_d\left(\ohat\right),\qquad
    \Lambda_d\left(\ohat\right):=\operatorname{diag}\left(W_mJ_m\,\ohat\cdot\nabla r_d\big|_m\right),                                                 \\
    \label{eq:general-boundary}
    \Gamma_f\left(\ohat\right) & \approx \operatorname{diag}\left(\omega^f_m J^f_m\,\ohat\cdot\hat{n}^f_m\right),\qquad f=1,\ldots,6,
  \end{align}
\end{subequations}
where $(S_{\ohat})_{\hat{m}m}\approx\int_e \ell_m\,\ohat\cdot\nabla\ell_{\hat{m}}\odif{V}$ is the streaming matrix, and, on face $f$, $\omega^f_m$ is the two-dimensional \gls{gll} weight (e.g., $\rho_j\rho_k$ on the faces $\mu=\pm1$), $J^f$ is the surface Jacobian (e.g., $\left|\pdv{\rvec}{\eta}\times\pdv{\rvec}{\xi}\right|$ on $\mu=\pm1$), and $\hat{n}^f$ is the outward unit normal.
The discrete form of \cref{eq:weak-form} on a single element is then
\begin{equation}
  \label{eq:sem-general}
  \sum_{f=1}^{6}E_f^{T}\,\Gamma_f\left(\ohat\right)\vec{\psi}^{*}_f
  -S_{\ohat}\vec{\psi}
  +\Sigma_t M\vec{\psi}
  =M\vec{q},
\end{equation}
where $E_f$ extracts the $(P+1)^2$ nodes of face $f$ from the volume vector, and $\vec{\psi}^{*}_f$ is the upwind trace on that face: the element's own angular flux at nodes where $\ohat\cdot\hat{n}^f>0$, and the angular flux from the neighboring element or the boundary condition where $\ohat\cdot\hat{n}^f<0$.
On a curved face, $\ohat\cdot\hat{n}^f$ can change sign along the face, so this choice is made node by node.

\Cref{eq:sem-general} is never assembled as a dense $(P+1)^3\times(P+1)^3$ element matrix.
Applying $S_{\ohat}$ to a vector requires only the three one-dimensional derivative sweeps $D_d$ and their transposes, each costing $\mathcal{O}(P^4)$ operations per element, together with $\mathcal{O}(P^3)$ diagonal scalings, compared with $\mathcal{O}(P^6)$ operations for a dense element matrix; only the nodal geometric factors need to be stored.
This \emph{sum factorization} is the matrix-free approach used by nekRS \cite{fischerNekRSGPUacceleratedSpectral2022}.

\subsubsection{Rectilinear Elements}
The tensor-product structure becomes explicit when the element is a rectangular box with side lengths $h_x$, $h_y$, and $h_z$, mapped affinely onto the reference cube:
\begin{equation}
  x = x_0+\frac{h_x}{2}\left(\mu+1\right),\quad
  y = y_0+\frac{h_y}{2}\left(\eta+1\right),\quad
  z = z_0+\frac{h_z}{2}\left(\xi+1\right).
\end{equation}
The geometric factors in \cref{eq:geometric-factors} are then constant: $J=\frac{h_xh_yh_z}{8}$, $\nabla\mu=\frac{2}{h_x}\hat{x}$, $\nabla\eta=\frac{2}{h_y}\hat{y}$, and $\nabla\xi=\frac{2}{h_z}\hat{z}$.
For example, $\Lambda_1=\ohat_x\frac{h_yh_z}{4}\operatorname{diag}\left(W_m\right)$, and since $\operatorname{diag}\left(W_m\right)=\hat{M}\otimes\hat{M}\otimes\hat{M}$ with $\hat{M}$ and $\hat{S}=\hat{D}^{T}\hat{M}$ as evaluated in \cref{eq:mass-gll-integral,eq:streaming-gll-integral} below, $D_1^{T}\Lambda_1=\ohat_x\frac{h_yh_z}{4}\left(\hat{M}\otimes\hat{M}\otimes\hat{S}\right)$.
Every term in \cref{eq:sem-general} collapses in this way into a Kronecker product of one-dimensional matrices, giving:
\begin{multline}
  \label{eq:sem-kronecker}
  \ohat_x\left(M_z\otimes M_y\otimes B_x\right)\vec{\psi}^{*}
  +\ohat_y\left(M_z\otimes B_y\otimes M_x\right)\vec{\psi}^{*}
  +\ohat_z\left(B_z\otimes M_y\otimes M_x\right)\vec{\psi}^{*}\\
  -\Big[\ohat_x\left(M_z\otimes M_y\otimes S_x\right)
    +\ohat_y\left(M_z\otimes S_y\otimes M_x\right)
  +\ohat_z\left(S_z\otimes M_y\otimes M_x\right)\Big]\vec{\psi}\\
  +\Sigma_t\left(M_z\otimes M_y\otimes M_x\right)\vec{\psi}
  =\left(M_z\otimes M_y\otimes M_x\right)\vec{q},
\end{multline}
where the $x$-direction matrix appears rightmost in each Kronecker product because the $x$ index varies fastest in the node ordering $m$.
The first three terms arise from the surface integral in \cref{eq:weak-form}, with $\vec{\psi}^{*}$ the upwind trace defined after \cref{eq:sem-general}.
Because the boundary matrices are nonzero only for the end nodes, only the face entries of $\vec{\psi}^{*}$ contribute, and since each face of a box has a constant normal, the upwind choice is made face by face.

Each one-dimensional matrix is a scaled copy of a reference matrix on $[-1,1]$:
\begin{equation}
  M_x=\frac{h_x}{2}\hat{M},\quad S_x=\hat{S},\quad B_x=\hat{B},
\end{equation}
and likewise in $y$ and $z$ with $h_y$ and $h_z$.
The streaming matrix carries no scale factor because the $\frac{h_x}{2}$ from $\odif{x}$ cancels the $\frac{2}{h_x}$ from $\pdv{}{x}$.
With rows indexed by the test function and columns by the trial function, for $\hat{\imath},i\in\{0,1,\ldots,P\}$:
\begin{subequations}
  \label{eq:kronecker-matrix-defs}
  \begin{align}
    \label{eq:mass-matrix}
    \left(\hat{M}\right)_{\hat{\imath}i} & := \int_{-1}^{1}\ell_{\hat{\imath}}(\mu)\,\ell_i(\mu)\odif{\mu}                \\
    \label{eq:streaming-matrix}
    \left(\hat{S}\right)_{\hat{\imath}i} & := \int_{-1}^{1}\odv{\ell_{\hat{\imath}}(\mu)}{\mu}\,\ell_i(\mu)\odif{\mu}     \\
    \label{eq:boundary-matrix}
    \left(\hat{B}\right)_{\hat{\imath}i} & := \Big[\ell_{\hat{\imath}}(\mu)\,\ell_i(\mu)\Big]_{-1}^{1}
    =\delta_{\hat{\imath}P}\delta_{iP}-\delta_{\hat{\imath}0}\delta_{i0},
  \end{align}
\end{subequations}
where the final equality in \cref{eq:boundary-matrix} holds because the endpoints $\mu=\pm1$ are \gls{gll} nodes.
Integrating \cref{eq:streaming-matrix} by parts shows that the three reference matrices are related by $\hat{B}=\hat{S}+\hat{S}^{T}$.

If we approximate the integrals in \cref{eq:mass-matrix,eq:streaming-matrix} using the \gls{gll} quadrature set of \cref{eq:gll}, noting that the nodes of the Lagrange interpolants are also the \gls{gll} nodes, so that $\ell_i(\mu_p)=\delta_{ip}$, we find:
\begin{subequations}
  \begin{equation}
    \label{eq:mass-gll-integral}
    \begin{split}
      \left(\hat{M}\right)_{\hat{\imath}i} & \approx \sum_{p=0}^{P}\rho_p\,\ell_{\hat{\imath}}(\mu_p)\,\ell_i(\mu_p) \\
                                           & =\sum_{p=0}^{P}\rho_p\,\delta_{\hat{\imath}p}\delta_{ip}               \\
                                           & =\rho_i\,\delta_{\hat{\imath}i}
    \end{split}
  \end{equation}
  \begin{equation}
    \label{eq:streaming-gll-integral}
    \begin{split}
      \left(\hat{S}\right)_{\hat{\imath}i} & = \sum_{p=0}^{P}\rho_p\,\odv{\ell_{\hat{\imath}}}{\mu}\bigg|_{\mu_p}\ell_i(\mu_p) \\
                                           & =\sum_{p=0}^{P}\rho_p\,\odv{\ell_{\hat{\imath}}}{\mu}\bigg|_{\mu_p}\delta_{ip}        \\
                                           & =\rho_i\,\hat{D}_{i\hat{\imath}},
    \end{split}
  \end{equation}
\end{subequations}
where $\hat{D}$ is the \gls{gll} differentiation matrix of \cref{eq:gll-diff}.
The integrand in \cref{eq:mass-matrix} is of degree $2P$, so \cref{eq:mass-gll-integral} is an approximation; the integrand in \cref{eq:streaming-matrix} is of degree $2P-1$, so \cref{eq:streaming-gll-integral} is exact.
In matrix form, $\hat{M}=\operatorname{diag}\left(\rho_0,\rho_1,\ldots,\rho_P\right)$ and $\hat{S}=\hat{D}^{T}\hat{M}$; that is, \cref{eq:mass-gll-integral} implies that the mass matrix defined in \cref{eq:mass-matrix} will be \emph{diagonal}.
The boundary matrix, \cref{eq:boundary-matrix}, requires no quadrature.
We can then write the entries of the Kronecker products in \cref{eq:sem-kronecker} using \cref{eq:mass-gll-integral,eq:streaming-gll-integral,eq:boundary-matrix}:
\begin{subequations}
  \begin{align}
    \left(M_z\otimes M_y\otimes M_x\right)_{\hat{m}m} & \approx \frac{h_xh_yh_z}{8}\,\rho_i\rho_j\rho_k\,\delta_{\hat{\imath}i}\delta_{\hat{\jmath}j}\delta_{\hat{k}k}                                \\
    \left(M_z\otimes M_y\otimes S_x\right)_{\hat{m}m} & \approx \frac{h_yh_z}{4}\,\rho_i\rho_j\rho_k\,\hat{D}_{i\hat{\imath}}\,\delta_{\hat{\jmath}j}\delta_{\hat{k}k}                               \\
    \left(M_z\otimes S_y\otimes M_x\right)_{\hat{m}m} & \approx \frac{h_xh_z}{4}\,\rho_i\rho_j\rho_k\,\delta_{\hat{\imath}i}\,\hat{D}_{j\hat{\jmath}}\,\delta_{\hat{k}k}                             \\
    \left(S_z\otimes M_y\otimes M_x\right)_{\hat{m}m} & \approx \frac{h_xh_y}{4}\,\rho_i\rho_j\rho_k\,\delta_{\hat{\imath}i}\delta_{\hat{\jmath}j}\,\hat{D}_{k\hat{k}}                               \\
    \left(M_z\otimes M_y\otimes B_x\right)_{\hat{m}m} & \approx \frac{h_yh_z}{4}\,\rho_j\rho_k\left(\delta_{\hat{\imath}P}\delta_{iP}-\delta_{\hat{\imath}0}\delta_{i0}\right)\delta_{\hat{\jmath}j}\delta_{\hat{k}k},
  \end{align}
\end{subequations}
with the $y$- and $z$-direction boundary terms following analogously.
The streaming and boundary entries are approximate only through the lumped mass factors in the transverse directions.

Although a reactor mesh is rarely made only of boxes, the rectilinear form remains useful in practice.
A Kronecker-product operator can be inverted cheaply by fast diagonalization, which applies the inverse through one-dimensional eigendecompositions in $\mathcal{O}(P^4)$ operations per element. % TODO cite Lynch, Rice, and Thomas (1964) for fast diagonalization
A curved element's operator, \cref{eq:sem-general}, does not factor this way, but replacing each element with a nearby box gives an inexpensive element-local preconditioner for it; Nek5000 and nekRS use this approach in their Schwarz smoothers \cite{fischerNekRSGPUacceleratedSpectral2022}.
